\scanpage{53}
% heading and introduction
\uline{Exposé II} \inserted{Cas d'action \deleted{\uline{\unsure{négligeable}}} du groupe fondamental} (Une variante du th. de monodromie) \deleted{\unsure{pour}}

Nous donnons dans le \deleted{\unsure{les articles|les autres}} présent exposé un résultat de nature analogue au théorème de monodromie de l'exposé précédent, \deleted{\unsure{pour}} à titre de nouvelle illustration de la méthode qui y était employée. Il ne servira plus dans la suite de ce séminaire.

\uline{1. Énoncé du résultat, et Corollaires.}

\deleted{\unsure{Énon}} Le résultat \deleted{essentiel} \inserted{principal du présent exposé} est le suivant.

% theorem, lines 7-12
\uline{Théorème} 1.1. Soient $k$ un corps \deleted{\unsure{transcendant} \unsure{quelconque}, tel que le \unsure{groupe} \unsure{soit} \unsure{de} \unsure{caract.} de \unsure{$k$}} \deleted{\unsure{sous-extension de type fini algébriquement}} de type fini (au sens absolu), $X$ un schéma sur \inserted{$S=\operatorname{Spec}k$} $k$, normal, localement de type fini, \deleted{\unsure{et}} géométriquement connexe, $\unsure{\xi|\zeta}$ un point géométrique \inserted{de $X$}, d'où (SGA 1 \uline{VIII}) une suite exacte
\[
1\longrightarrow\pi_1(\overline X,\xi)\longrightarrow\pi_1(X,\xi)\longrightarrow\pi_1(S,\xi)\longrightarrow1
\tag*{(1.1.1)}
\]
% lines 14-18
avec \deleted{\unsure{$\overline k$}} $\overline X=X\otimes_k\overline k$, $\overline k$ étant la clôture séparable de $k$ dans $k(\xi)$, (et $\xi$ désignant aussi les points géométriques de $S$ resp. $\overline X$ déduits canoniquement du \deleted{\unsure{$\xi$}} point géométrique $\xi$ de $X$). \deleted{Soit \unsure{$I$} \unsure{le} \unsure{groupe} \unsure{Considérons le plus grand}}

% lines 19-23
\deleted{\unsure{que}} Soient $\ell$ un nombre premier distinct de la caractéristique de $k$, \deleted{\unsure{et}} $\pi$ un sous-groupe ouvert de $\pi_1(X,\xi)$, $\pi_0$ son image dans $\pi_1(S,\xi)$, $\overline I=\pi\cap\pi_1(\overline X,\xi)$, de sorte qu'on a une suite exacte
\[
1\mathrel{\text{---}}\overline I\mathrel{\text{---}}\pi\mathrel{\text{---}}\pi_0\longrightarrow1.
\tag*{(1.1.2)}
\]
% lines 24-27
Soit $I_{\deleted{\unsure{\ell}}}(\ell)$ la composante $\ell$-primaire du plus grand quotient profini abélien $I_{\mathrm{ab}}$ de $I$, \deleted{Alors} d'où \inserted{par passage au quotient} une extension
\[
1\mathrel{\text{---}}I(\ell)\mathrel{\text{---}}\pi'_{\deleted{\unsure{\ell}}}\longrightarrow\pi_0\mathrel{\text{---}}1.
\tag*{(1.1.3)}
\]
% lines 28-32
Alors \deleted{$\unsure{[\pi',I(\ell)]}$ \unsure{est un groupe}} $I(\ell)$ est un module de type fini sur $\mathbb{Z}_\ell$, et $\overline{[\pi',I(\ell)]}$ \deleted{est de} d'indice fini dans $I(\ell)$, en d'autres termes \deleted{\unsure{$I(\ell)$}} le plus grand quotient profini de $I(\ell)$ sur lequel $\pi_0$ opère trivialement est un groupe \uline{fini}.
% archive sheet number: 52
